\section{本章小结}
\label{sec:rvc-summary}

本章从两台传感器的问题出发：两台同时报偏高的概率、一台读数已知时另一台的规律、平均以后的波动，都要求知道几个量在同一次观测中怎样共同取值。把定义在同一样本空间上的随机变量排成随机向量，这份共同规律就是联合分布。离散时它是一张概率表，连续时是单位面积上的质量即联合密度，两类模型都可以用联合分布函数描述，矩形概率由四个角上的值确定。

从联合分布出发，本章得到的几种对象可以放在同一张概率表上看。表~\ref{tab:rvc-summary}按传感器表列出它们的算法与结果。

\begin{BookInlineTable}
  \centering
  \caption{从联合分布得到的对象（传感器表）}
  \label{tab:rvc-summary}
  \begin{tabular}{lll}
    \toprule
    对象 & 在概率表上的操作 & 结果 \\
    \midrule
    边缘分布 & 按行、按列求和 & $(0.3,0.4,0.3)$ \\
    独立性 & 每格与边缘之积比较 & 不独立 \\
    均值向量 & 按质量加权平均 & $(1,1)^{\top}$ \\
    协方差 & 中心化乘积加权求和 & $0.4$ \\
    条件分布 & 取出一行，除以行和 & 如$(\tfrac23,\tfrac13,0)$ \\
    条件均值 & 对每行求均值 & $\tfrac13+\tfrac23x$ \\
    \bottomrule
  \end{tabular}
\end{BookInlineTable}

边缘化把其余分量的信息汇总掉，不能逆转：相同的边缘可以对应许多不同的联合分布。只有在独立时，联合分布才由边缘拼成乘积；判别时，概率表要求每个格子分解，密度要求边缘乘积本身就是联合密度。独立带来两项计算上的便利：乘积的期望可以分开计算（两个期望有限时），条件分布不随条件改变。独立同分布同时要求同分布与相互独立，两项不能互相替代；两两独立也不足以推出相互独立。

均值向量是质量分布的重心，协方差矩阵描述质量沿各个方向的散开程度。它对称半正定，$\bm{a}^{\top}\bm{\Sigma}\bm{a}$是$\bm{a}^{\top}\bm{X}$的方差，奇异时质量压在一个超平面上。有限和的方差展开中，交叉项来自分量之间的协方差，它决定平均能把波动降低多少：各读数方差同为$\sigma^{2}$、两两相关系数同为$\rho>0$时，无论平均多少个读数，方差都降不到$\rho\sigma^{2}$以下。相关系数介于$-1$与$1$之间，衡量向一条直线集中的程度；独立推出不相关，反之不成立，零相关只排除线性关联。

条件分布在概率表中是取出一行再重新归一，在密度中是取一条切片再除以$f_X(x)$，只在分母为正时使用。条件均值是组内平均，全期望公式说总平均是组平均的加权平均，全方差公式把总方差分成组内方差的平均与组平均之间的方差。单个零概率条件下的条件密度不由联合分布唯一确定，严格的处理在第~\ref{chap:conditional-distributions}章完成。

下一章用这些结构考察具体的多元模型：多项计数的和为常数，造成奇异协方差；多元高斯分布由均值向量和协方差矩阵完全确定；共同尺度可以让不相关的分量仍然相互依赖。按给定联合分布实际产生样本以及效率问题，在第~\ref{chap:sample-generation}章展开。
